\documentclass[11pt,a4paper]{article}
\usepackage[T1]{fontenc}
\usepackage[utf8]{inputenc}
\usepackage{lmodern,amsmath,amssymb}
\usepackage[margin=25mm]{geometry}
\usepackage[hidelinks]{hyperref}
\hypersetup{pdftitle={Why the abelian theory has no gap, in the language of the target box: its coupling does no},pdfauthor={navstokgap research working notes}}
\newcommand{\tightlist}{\setlength{\itemsep}{0pt}\setlength{\parskip}{0pt}}
\setlength{\emergencystretch}{3em}
\setcounter{secnumdepth}{0}
\begin{document}
\hypertarget{why-the-abelian-theory-has-no-gap-in-the-language-of-the-target-box-its-coupling-does-not-run}{%
\section{Why the abelian theory has no gap, in the language of the
target box: its coupling does not
run}\label{why-the-abelian-theory-has-no-gap-in-the-language-of-the-target-box-its-coupling-does-not-run}}

The founding question of this series, why the commutative theory is
gapless while the non-abelian one is expected to be gapped, has a
precise form once the strong-coupling gap is a box rather than a point
at infinity. The continuous-time expansion of
\href{kogut-susskind-strong-coupling-explicit.md}{the Kogut--Susskind
note} is group-blind: it uses only that the electric term is diagonal
with a least nonzero Casimir, that the plaquette term is bounded and
acts on four links, and Gauss's law. For compact \(U(1)\) it gives a gap
\(\Delta\ge4\lambda\) with \(\varepsilon=g^2/2\), for \(g^2\) above a
threshold of the same size as for \(SU(3)\), uniformly in the volume. So
\textbf{both theories are gapped inside the box.} What differs is
whether the flow from weak coupling reaches it. For \(SU(3)\) the
coupling runs,
\[\frac{d\,g^2}{d\log(a)}=2b_0g^4+O(g^6),\qquad b_0=\frac{11}{16\pi^2}>0,\]
and the conjecture is that it keeps running until it enters the box. For
\(U(1)\) the pure gauge theory has \(b_0=0\): the coupling does not run,
the weak-coupling theory is the free photon at every scale, and it never
reaches the box; the massless Coulomb phase of Guth (Phys. Rev.~D 21
(1980) 2291) and Fröhlich--Spencer (Commun. Math. Phys. 83 (1982) 411;
both metadata level) is the rigorous form of this statement, and the
gapped set of the lattice \(U(1)\) theory is \((g_c,\infty)\) with
\(g_c>0\). In the language of \href{gapped-set-critical-coupling.md}{the
critical-coupling note}, T2\('\) for \(SU(3)\) is the statement that the
renormalization flow has no fixed point between the ultraviolet and the
box, and for \(U(1)\) the whole weak-coupling region is a line of fixed
points. Constants explicit; nothing promoted.

\hypertarget{the-strong-coupling-box-does-not-see-the-group}{%
\subsection{1. The strong-coupling box does not see the
group}\label{the-strong-coupling-box-does-not-see-the-group}}

Run the expansion of the Kogut--Susskind note for compact \(U(1)\) with
the Wilson term \((2/g^2)\sum_p(1-\cos\theta_p)\). The electric term is
\((g^2/2)\sum_\ell n_\ell^2\), diagonal in the charge basis, with least
nonzero eigenvalue \(\varepsilon=g^2/2\) per link. The plaquette
operator \(w_p=e^{i\theta_p}+e^{-i\theta_p}\) has norm \(2\) and shifts
the charge on the four links of \(p\) by \(\pm1\). Gauss's law,
\(\sum_{\ell\ni x}\pm n_\ell=0\) at every site, makes every excited set
a union of closed charged loops, so \(|S_k|\ge4\). The per-step factor
of the adjacent-growth count is \(u=128e/(g^2(\varepsilon-\lambda))\)
with \(N\to1\), and the rest is verbatim. The abelian theory is
therefore gapped, uniformly in the volume, for \(g^2\) above an explicit
threshold of the same order as for \(SU(3)\), with a gap approaching the
flux-loop energy \(4\cdot\frac{g^2}2=2g^2\) in units \(\hbar c/a\). This
is the confining phase of compact \(U(1)\), and it is the same mechanism
as for \(SU(3)\): electric flux costs energy per link, and Gauss's law
forces flux to close.

\hypertarget{what-differs-is-the-flow}{%
\subsection{2. What differs is the
flow}\label{what-differs-is-the-flow}}

Let \(g(a)\) be the effective coupling of the lattice theory at spacing
\(a\), defined by any renormalization scheme that is an exact low-energy
reduction. The target box of \href{strong-coupling-target-box.md}{the
target-box note} is the region \(g^2\gtrsim4\times10^2\) with few-link
corrections of controlled local norm, and the Hamiltonian at any scale
that lies in it has a gap.

\emph{\(SU(3)\).} Asymptotic freedom gives, at weak coupling,
\(g^2(2a)=g^2(a)+2b_0\log2\,g^4(a)+O(g^6)\) with \(2b_0\log2=0.0966\),
so the coupling grows as the scale coarsens. The mass-gap conjecture is
the statement that this growth continues, with corrections staying in
the class \(\mathcal C\), until \(g^2\) enters the box; equivalently,
that the flow has no fixed point at intermediate coupling. The
intermediate region of \href{mass-gap-position.md}{the position note} §6
is the stretch of the flow in which neither the perturbative form of the
running nor the box applies.

\emph{\(U(1)\).} The pure abelian theory has no self-interaction of the
gauge field, so \(b_0=0\) and, at weak coupling, \(g(2a)=g(a)\) to all
orders: the free photon is a fixed point of the flow at every value of
the coupling. The weak-coupling theory therefore never moves toward the
box. Guth and Fröhlich--Spencer prove that for \(g\) small the lattice
theory is in a massless phase with power-law decay of the Wilson loop
and a massless photon, and the strong-coupling expansion proves a gap
for \(g\) large; the gapped set is \((g_c,\infty)\) with
\(0<g_c<\infty\), and the transition at \(g_c\) is the zero-temperature
phase transition that T2\('\) forbids for \(SU(3)\).

\hypertarget{the-three-earlier-isolations-of-the-non-abelian-structure-in-this-language}{%
\subsection{3. The three earlier isolations of the non-abelian
structure, in this
language}\label{the-three-earlier-isolations-of-the-non-abelian-structure-in-this-language}}

\href{mass-gap-position.md}{The position note} §4 lists three places
where the non-abelian structure enters: the commutator potential of the
zero-momentum sector, the absence of a gauge-invariant operator linear
in the electric field, and the one-loop valley potential. All three are
statements about the first coefficient of the running: the commutator is
what makes \(b_0\ne0\), the valley potential is its finite-volume image,
and the absence of a linear gauge-invariant operator is why the photon
channel through which the abelian gap closes does not exist. In the box
language they are the reasons the \(SU(3)\) flow leaves the ultraviolet
at all; the open problem is whether it arrives.

\hypertarget{what-this-settles-and-what-it-leaves}{%
\subsection{4. What this settles and what it
leaves}\label{what-this-settles-and-what-it-leaves}}

\emph{Settled.} The gap is a property of the box, and the box is
group-blind: confinement at strong coupling is the same for \(U(1)\) and
\(SU(3)\). The difference between the theories is the flow into the box,
and for \(U(1)\) the flow is absent. ``Commutative fields have no gap''
means: their coupling does not run, so the weak-coupling theory is a
fixed point outside the box.

\emph{Left.} For \(SU(3)\) the flow leaves the ultraviolet with
\(b_0>0\) and the box is at \(g^2\sim4\times10^2\). Whether the flow
arrives is the intermediate-region problem, and every quantity in it is
now explicit: the starting slope \(2b_0\log2=0.0966\) per doubling, the
class \(\mathcal C\) that has to be preserved, and the box's edge and
tolerance.

\hypertarget{consequence-for-state}{%
\subsection{5. Consequence for STATE}\label{consequence-for-state}}

The founding question has its answer at the level of the current map:
the strong-coupling gap is group-blind, the flow is what distinguishes
the groups, and the abelian theory's coupling does not run. The
non-abelian mass gap is the statement that the running, which begins
with \(b_0=11/(16\pi^2)\) for \(SU(3)\), does not stop before the box.
\end{document}
